%=====================================================================
\chapter{Mixed Methods and Triangulation}\label{ch:mixed}
%=====================================================================
\locator{3}{Part~3 --- Research Designs for Computing}

\begin{outcomes}
By the end of this chapter you will be able to:
\begin{tight}
  \item \textbf{State} a reason for mixing that is specific to your question.
  \item \textbf{Select} among convergent, explanatory sequential, exploratory
    sequential and embedded designs.
  \item \textbf{Notate} a design so a reader knows the sequence and the
    priority.
  \item \textbf{Integrate} strands with a joint display rather than
    juxtaposing them.
  \item \textbf{Recognise} when mixing adds nothing and should be dropped.
\end{tight}
\end{outcomes}

\begin{prereq}
\chref{philosophy} for pragmatism, \chref{experiments} and \chref{survey} for
the quantitative strands, \chref{qualitative} and \chref{casestudy} for the
qualitative ones.
\end{prereq}

\begin{keyterms}
mixed methods $\bullet$ convergent design $\bullet$ explanatory sequential
$\bullet$ exploratory sequential $\bullet$ embedded design $\bullet$ priority
$\bullet$ integration $\bullet$ joint display $\bullet$ meta-inference
$\bullet$ triangulation $\bullet$ divergence
\end{keyterms}

\section{Mixing Needs a Reason}

``We conducted a survey and also some interviews'' is not mixed methods
research. It is two studies in one document. Mixed methods means the strands
are \emph{integrated} --- one informs, explains, or is explained by the other,
and the combined inference is something neither could have produced alone.

\begin{center}
\small
\begin{tabular}{@{}L{4.0cm}L{10.3cm}@{}}
\toprule
\textbf{Legitimate reason} & \textbf{What it looks like} \\
\midrule
Explanation
  & The quantitative result is puzzling; qualitative work explains the
    mechanism \\
Development
  & Qualitative work identifies the constructs; a quantitative instrument
    measures them at scale \\
Triangulation
  & Two methods addressing the same question, where agreement increases
    confidence and disagreement is informative \\
Complementarity
  & Different facets of one phenomenon: how often it happens, and what it is
    like \\
Initiation
  & Deliberately seeking contradiction, to provoke reframing \\
\bottomrule
\end{tabular}
\end{center}

\begin{alertbox}[If you cannot name the reason, do not mix]
Mixed methods costs roughly twice as much as a single-method study and requires
competence in both traditions. A graduate scholar with limited time who mixes
without a reason produces two mediocre strands instead of one good one.

\medskip
The test: write the sentence ``the qualitative strand is necessary because the
quantitative strand cannot \underline{\hspace{25mm}}''. If you cannot complete
it specifically, drop a strand and do the remaining one properly.
\end{alertbox}

\section{Four Designs and Their Notation}

Notation makes the design unambiguous. Capitals indicate the dominant strand,
a plus sign indicates concurrent strands, an arrow indicates sequence.

\begin{center}
\small
\begin{longtable}{@{}L{3.2cm}L{2.3cm}L{4.0cm}L{4.2cm}@{}}
\toprule
\textbf{Design} & \textbf{Notation} & \textbf{Sequence} & \textbf{Use when} \\
\midrule
\endfirsthead
\toprule
\textbf{Design} & \textbf{Notation} & \textbf{Sequence} & \textbf{Use when} \\
\midrule
\endhead
\bottomrule
\endfoot
Convergent parallel
  & QUAN + QUAL
  & Both at once, merged at interpretation
  & You want corroboration or a fuller picture, and neither strand needs the
    other's results first \\
Explanatory sequential
  & QUAN $\rightarrow$ qual
  & Quantitative first, qualitative to explain it
  & A result is surprising, or you need the mechanism behind an effect. The
    most common design in computing \\
Exploratory sequential
  & QUAL $\rightarrow$ quan
  & Qualitative first, then measure at scale
  & The constructs are unknown; you must find out what to measure before you
    can measure it \\
Embedded
  & QUAN(qual)
  & A small strand inside a larger design
  & Interviews inside an experiment to understand why participants behaved as
    they did \\
\end{longtable}
\end{center}

\armfig{
\begin{tikzpicture}[
  qn/.style={rounded corners=2pt, draw=tradEMPc!70, fill=tradEMPc!10,
             text width=2.0cm, align=center, font=\scriptsize,
             minimum height=0.85cm},
  ql/.style={rounded corners=2pt, draw=tradQUALc!70, fill=tradQUALc!10,
             text width=2.0cm, align=center, font=\scriptsize,
             minimum height=0.85cm},
  intg/.style={rounded corners=2pt, draw=goldhl!80, fill=goldhl!12,
             text width=1.9cm, align=center, font=\scriptsize,
             minimum height=0.85cm},
  lb/.style={font=\scriptsize\bfseries, text=primary},
  ar/.style={-{Stealth[length=2mm]}, draw=secondary, semithick}]

  \node[lb] at (-0.2,1.35) {Convergent};
  \node[qn] (c1) at (0.9,0.55) {QUAN};
  \node[ql] (c2) at (0.9,-0.45) {QUAL};
  \node[intg] (c3) at (3.4,0.05) {merge and interpret};
  \draw[ar] (c1) -- (c3); \draw[ar] (c2) -- (c3);

  \node[lb] at (5.9,1.35) {Explanatory sequential};
  \node[qn] (e1) at (6.3,0.05) {QUAN};
  \node[ql] (e2) at (8.6,0.05) {qual};
  \node[intg] (e3) at (10.9,0.05) {explain the result};
  \draw[ar] (e1) -- node[above, font=\scriptsize, text=neutral] {select cases} (e2);
  \draw[ar] (e2) -- (e3);

  \node[lb] at (-0.2,-1.65) {Exploratory sequential};
  \node[ql] (x1) at (0.9,-2.45) {QUAL};
  \node[qn] (x2) at (3.4,-2.45) {quan};
  \node[intg] (x3) at (5.9,-2.45) {generalise};
  \draw[ar] (x1) -- node[above, font=\scriptsize, text=neutral] {build instrument} (x2);
  \draw[ar] (x2) -- (x3);

  \node[lb] at (8.9,-1.65) {Embedded};
  \node[qn, minimum height=1.5cm, text width=3.0cm] (b1) at (9.9,-2.45) {\\ QUAN};
  \node[ql, text width=1.5cm, minimum height=0.6cm] (b2) at (9.9,-2.75) {qual};
\end{tikzpicture}}
{Four mixed methods designs. The gold box is the integration point --- the part
that distinguishes mixed methods from two studies bound together, and the part
most often missing.}{mixeddesigns}

\begin{worked}[An explanatory sequential design]
\textbf{Strand 1 (QUAN).} A survey of 280 developers finds that teams using a
mandatory code review checklist report \emph{lower} satisfaction with review
quality than teams without one --- the opposite of the prediction.

\medskip
\textbf{The integration decision.} The surprising result determines who is
interviewed: not a fresh sample, but a purposive selection from the survey
respondents. Specifically, the extreme cases --- checklist users with the
lowest satisfaction scores --- plus a contrasting group of checklist users with
high scores. This selection \emph{is} the integration; it could not have been
made without the quantitative result.

\medskip
\textbf{Strand 2 (qual).} Twelve interviews. The explanation that emerges:
where review is time-boxed below the time the checklist takes to complete,
reviewers satisfy the checklist mechanically and stop looking for anything
else. Satisfaction is low because reviewers know they are performing
compliance rather than reviewing.

\medskip
\textbf{Meta-inference.} Neither strand supports this alone. The survey
establishes the association across 280 people; the interviews supply the
mechanism and the scope condition. The combined claim --- checklists reduce
perceived review quality \emph{when review time is constrained below checklist
completion time} --- is testable, and \chref{experiments} could now test it.

\medskip
All figures illustrative.
\end{worked}

\section{Integration Is the Hard Part}

\begin{center}
\small
\begin{tabular}{@{}L{3.2cm}L{11.1cm}@{}}
\toprule
\textbf{Where} & \textbf{How} \\
\midrule
Design level   & One strand's sampling or instrument is determined by the
                 other's results \\
Method level   & Building the questionnaire from interview codes; selecting
                 interviewees from survey responses \\
Analysis level & Quantitising qualitative data (counting code occurrences) or
                 qualitising quantitative data (creating profiles from
                 clusters) \\
Reporting level& The joint display, and the meta-inference paragraph \\
\bottomrule
\end{tabular}
\end{center}

\subsection{The joint display}

A \term{joint display} is a table or figure placing the two strands side by
side so the reader can see the integration rather than being told about it.

\begin{center}
\small
\begin{tabular}{@{}L{2.9cm}L{3.5cm}L{3.5cm}L{3.6cm}@{}}
\toprule
\textbf{Finding} & \textbf{Quantitative} & \textbf{Qualitative} & \textbf{Meta-inference} \\
\midrule
Checklist and satisfaction
  & Mean satisfaction 3.1 vs 3.8, $d = -0.41$
  & Reviewers describe ``ticking boxes'' under time pressure
  & Confirms and explains: the effect is mediated by time-boxing \\
Experience
  & No moderation by years of experience
  & Senior reviewers describe ignoring the checklist entirely
  & \textbf{Divergence}: the null moderation may mask two opposite behaviours
    cancelling out. Warrants re-analysis \\
\bottomrule
\end{tabular}
\end{center}

\begin{conceptbox}[Divergence is the most valuable row in the table]
When the strands disagree, resist the urge to explain it away. Four things it
can mean, and they are distinguishable:

\begin{tightnum}
  \item \textbf{The construct differs.} The survey and the interviews were
    about subtly different things.
  \item \textbf{Different populations.} Respondents and interviewees were not
    comparable.
  \item \textbf{One strand is wrong.} An underpowered quantitative result, or a
    qualitative interpretation over-fitted to a vivid case.
  \item \textbf{Genuine complexity.} Two mechanisms operate in opposite
    directions, and the aggregate hides both --- as in the experience row
    above.
\end{tightnum}

\medskip
The fourth is the interesting case, and it is the reason mixing was worth
doing. Report the divergence and your reasoning about which of the four it is.
\end{conceptbox}

\begin{pitfall}[Four ways mixed methods goes wrong]
\begin{tight}
  \item \textbf{No integration.} Chapter 4 quantitative, Chapter 5 qualitative,
    a conclusion mentioning both. This is the default failure.
  \item \textbf{Unstated priority.} The reader cannot tell which strand carries
    the claim, so cannot judge whether its evidence suffices.
  \item \textbf{Sequence violated.} An exploratory sequential design in which
    the questionnaire was actually written before the interviews, which then
    served to justify it.
  \item \textbf{Incompatible paradigms unacknowledged.} Treating interview
    counts as a measure of prevalence, which imports statistical logic into
    purposively sampled data. Counting codes is legitimate as a descriptive
    device within the sample; it is not an estimate of anything.
\end{tight}
\end{pitfall}

\begin{traditions}{mixing methods}
\tradEMP{Values the qualitative strand for explaining a mechanism and for
generating hypotheses. Insists that the quantitative strand's inference is not
weakened by the qualitative one's sampling.}
\tradDSR{Naturally mixed: qualitative problem characterisation, artefact
construction, quantitative evaluation, qualitative study of use. FEDS
(\chref{designscience}) is effectively a mixed methods evaluation plan.}
\tradQUAL{Cautions that integration can subordinate qualitative work to a
supporting role --- ``quotes to illustrate the numbers''. Priority notation
exists partly to make that subordination visible when it happens.}
\tradFORM{Little direct role; but formal analysis plus empirical measurement of
whether the modelled effect matters at realistic scale is a mixing of a
different kind.}
\end{traditions}

\begin{lab}[Design a mixed study, or decide not to]
\begin{tightnum}
  \item Complete the sentence: ``the second strand is necessary because the
    first cannot \underline{\hspace{25mm}}''. If you cannot, stop here and note
    that as your finding.
  \item Choose a design and write its notation, including which strand
    dominates.
  \item Specify the integration point precisely: what exactly does one strand
    determine about the other?
  \item Draft an empty joint display with your expected finding rows.
  \item Write what you would do if the strands diverge, using the four
    categories above --- \emph{before} you have any data.
\end{tightnum}
\end{lab}

\begin{checkpoint}
\begin{tightnum}
  \item Why is ``we did a survey and interviews'' not a mixed methods study?
  \item You find a counter-intuitive quantitative result. Which design follows,
    and how does the first strand determine the second's sampling?
  \item What does QUAL $\rightarrow$ quan tell a reader that ``mixed methods''
    does not?
  \item Your strands disagree. Give the four possible explanations and one
    piece of evidence that would discriminate among them.
\end{tightnum}
\end{checkpoint}

\begin{chaptersummary}
\begin{tight}
  \item Mixed methods means integrated strands, not two studies in one
    document. Name the reason for mixing or drop a strand.
  \item Four designs: convergent, explanatory sequential, exploratory
    sequential, embedded. Notate them with capitals for priority and arrows for
    sequence.
  \item Integration happens at design, method, analysis or reporting level.
    Design-level integration --- one strand determining the other's sampling
    --- is the strongest.
  \item The joint display shows the integration instead of asserting it.
  \item Divergence between strands is the most valuable finding, and there are
    four distinguishable reasons for it.
  \item Counting codes describes your sample; it does not estimate prevalence.
\end{tight}
\end{chaptersummary}

\begin{reviewq}
  \item Pragmatism is offered as the paradigm that licenses mixing. Is it a
    principled position or a way of avoiding the incompatibility problem?
    Argue both, with reference to \chref{philosophy}.
  \item Mixed methods doubles the cost. For a two-semester MS thesis, when is
    that ever the right allocation? Give a concrete example and say what would
    be sacrificed.
  \item Priority notation makes subordination of the qualitative strand
    visible. Does making it visible make it acceptable? Discuss with reference
    to a QUAN(qual) design.
  \item Design a study of your own topic in which the strands are genuinely
    likely to diverge, and say in advance what each of the four explanations
    would imply for your conclusion.
\end{reviewq}
